In generational populations of LLM agents, a name does at least two separable things. As a
stable handle, it lets agents find the same partner again, and any stable handle does that job
equally well. As something parents are said to give, it frames the population as a family:
either the parental act or merely the wording raised couple formation, while, in exploratory
analyses, a world-assigned ID lowered it below no handle at all. Only names that parents composed themselves raised investment in
newborns. Arbitrary group labels reproduced the minimal-group pattern, yet strengthened rather than
weakened coordination when added to individual handles. Two registered human parallels failed:
namesaking did not raise the namesake parent's investment, and the agents almost entirely
free-rode on a shared pot. One unregistered sentence about namesakes halved couple formation.
Together these results suggest that how agents are named, and how the naming is described, should
be treated as an experimental variable in multi-agent LLM research rather than as scaffolding.
